% \begin{figure*}[t]
% \centering
% \begin{promptbox}[System Message]
% You generate hidden reasoning traces for user-simulation SFT data.

% Given [HUMAN]'s history, [HUMAN]'s persona, the current
% conversation context, and [HUMAN]'s actual next response,
% explain the reasoning process that could lead [HUMAN]
% to write that response.
% \end{promptbox}
% % \vspace{4pt}

% \begin{promptbox}[User Message]
% <|Target User Evidence|>
% (reserved history block (*@$h$@*))
% <|End Target User Evidence|>

% <|Persona|>
% (induced persona (*@$\rho$@*))
% <|End Persona|>

% <|Interaction Context|>
% (current context (*@$x$@*))
% <|End Interaction Context|>

% <|Given Response|>
% (ground-truth response (*@$y$@*))
% <|End Given Response|>

% Generate the hidden reasoning trace that could lead
% [HUMAN] to write the given response.

% The trace should reason carefully about:
% - what [OTHER] or [OTHER - OP] most recently said
% - what exact point [HUMAN]'s response should target
% - [HUMAN]'s likely stance, emotion, goal, and style
% - plausible response length and level of detail

% Bring in supporting evidence from [HUMAN]'s history messages or persona.

% Do not reproduce the given response. Do not copy, quote, closely paraphrase, or reveal its final wording. The trace may describe the response's intent, stance, target, and style, but must not restate the response itself.

% Output only the reasoning trace text, with no XML tags or headings.
% \end{promptbox}


% \vspace{-6pt}

% \caption{
% Thinking trace elicitation prompt. A teacher model receives the user's history $h$, persona $\rho$,
% current context $x$, and ground-truth response $y$, and generates
% a reasoning trace explaining what could lead the user to write
% that response without reproducing it.
% }
% \label{fig:thinking_trace_prompt}
% \end{figure*}

\begin{figure*}[t]
\centering
\begin{adjustbox}{max width=\linewidth, scale=0.85}
\begin{minipage}{\linewidth}
\begin{promptbox}[System Message]
You generate hidden reasoning traces for user-simulation SFT data.
Given [HUMAN]'s history, [HUMAN]'s persona, the current
conversation context, and [HUMAN]'s actual next response,
explain the reasoning process that could lead [HUMAN]
to write that response.
\end{promptbox}
\end{minipage}
\end{adjustbox}
\ 

\begin{adjustbox}{max width=\linewidth, scale=0.85}
\begin{minipage}{\linewidth}
\begin{promptbox}[User Message]
<|Target User Evidence|>
(reserved history block (*@$h$@*))
<|End Target User Evidence|>
<|Persona|>
(induced persona (*@$\rho$@*))
<|End Persona|>
<|Interaction Context|>
(current context (*@$x$@*))
<|End Interaction Context|>
<|Given Response|>
(ground-truth response (*@$y$@*))
<|End Given Response|>
Generate the hidden reasoning trace that could lead
[HUMAN] to write the given response.
The trace should reason carefully about:
- what [OTHER] or [OTHER - OP] most recently said
- what exact point [HUMAN]'s response should target
- [HUMAN]'s likely stance, emotion, goal, and style
- plausible response length and level of detail
Bring in supporting evidence from [HUMAN]'s history messages or persona.
Do not reproduce the given response. Do not copy, quote, closely paraphrase, or reveal its final wording. The trace may describe the response's intent, stance, target, and style, but must not restate the response itself.
Output only the reasoning trace text, with no XML tags or headings.
\end{promptbox}
\end{minipage}
\end{adjustbox}
\vspace{-6pt}
\caption{
Thinking trace elicitation prompt. A teacher model receives the user's history $h$, persona $\rho$,
current context $x$, and ground-truth response $y$, and generates
a reasoning trace explaining what could lead the user to write
that response without reproducing it.
}
\label{fig:thinking_trace_prompt}
\end{figure*}